\begin{abstract}
DevSecOps evaluations frequently emphasize scanner counts, even though delivery decisions depend on whether findings survive remediation and violate explicit policy gates. SecFlowOps addresses this measurement gap with a reproducible evaluation framework that executes Trivy, Semgrep, Gitleaks, npm audit, OWASP ZAP, and OPA; normalizes findings; applies bounded deterministic and npm-based remediations; rescans patched workspaces; and records residual risk, policy outcomes, and runtime cost. The evaluation comprises a 432-run controlled protocol, a 90-run controlled size study, a 75-run open-source campaign with static adjudication, a NodeGoat npm-remediation study, an injected external recall study, and a targeted ZAP DAST probe. SecFlowOps reaches recall 1.0 on injected external ground truth and removes all injected target findings after remediation. In the 432-run protocol, mean residual findings decrease from 28.0 under scan-only configurations to 16.0 under the remediation configurations. On NodeGoat, npm-only remediation reduces findings from 78 to 32, while SecFlowOps reduces them to 31 by also removing a repository secret; residual critical findings remain blocked by policy. The results support SecFlowOps as an executable framework for auditing scanner-to-policy trajectories, not as proof that arbitrary projects can be automatically repaired or safely deployed.
\end{abstract}

\begin{IEEEkeywords}
DevSecOps, software security, vulnerability scanning, automated remediation, policy as code, reproducibility.
\end{IEEEkeywords}
